% Transcription — fonds Alexandre Grothendieck, shelfmark 85, batch 3 (pages 41-60).
% Established from the facsimile archives/batches/85/batch-03.pdf (a copy of
% 85.pdf fetched through the project's relay,
% https://grothendieck-relay.onrender.com/source/85.pdf), per the skill
% transcribe-grothendieck. The folder is ninety-seven pages in five batches;
% this is the third.
%
% Pass: Opus 5.5 (claude-opus-5-5), 2026-09-26 — first pass, unchecked against the pages by a human.
%
% Pages skipped: none as unrelated. Pages 41 and 57 are covers on pink card,
% each inscribed in ink at the top right and counted in pencil; they take no
% \page marker and their words become the section headings, with a \note.
% Page 41: « Systèmes de Schwartz dans A4, S4, A5, et classification des
% hyperpolygônes réguliers en car. zéro » « (15 p) » — exactly pages 42-56.
% Page 57: « Hyperpolygônes réguliers. » « (16 p) » — pages 58-73, of which
% only 58-60 fall in this batch. Pages transcribed: 42-56, 58-60 (eighteen).
% Pages summarised: none; where a crowded cancelled block (49, 55, 56) will
% not read, one \struck{\ill{}} or a \note stands for it.
%
% His pagination: none of his own on these leaves. The runs are followed by
% the pencil counts on the covers and by his numbered steps: « 2. », « 3) Cas
% du A4 », « 3. » (pp. 44-47); the lettered statements a)-f) of the
% Proposition on the tetrahedron (pp. 49-50) and a)-d) on the octahedron
% (pp. 52-53), with « Dém. de d) » and cases 1)-3) (pp. 54-56); I) and II)
% on pp. 58-59. No dates anywhere in the batch.
%
% What the batch is about. Pages 42-56: « systèmes de Schwartz » (his
% spelling) — triples (rho, rho', rho'') with rho rho' rho'' = 1 generating a
% finite group H — in A4, S4, A5. Page 42 tabulates the numerical types
% (orders) for the three groups, with counts 4, 9, 32 (total 45, or 16 modulo
% S5 for A5) and asks whether a Schwartz system is known up to conjugation
% once its terms are. Pages 43-44: finite subgroups of the automorphism group
% scheme of a conic X in characteristic 0; forms of A4 and S4, the
% equivalences F(S4) ≃ F(A4) ≃ (Ens)_4, Rev_4(S) ≃ Tors(S_4S) ≃ oriented
% cubes on S; rigidity of (G, rho, rho', rho''). Page 45: the triple of orders
% as first invariant; cases I, II, III (<= 3, 4, 5) for A4, S4, A5; start of
% the theorem for an oriented tetrahedron (rho_s of order 3). Page 46 (an
% inserted leaf, the argument of 45 resuming on 47): the S3-action on
% Tr(H) and on Schwartz systems; what « isomorphic » means in each case.
% Page 47: for a regular orientable combinatorial polyhedron with
% H = Aut^+(Pi), an element rho_phi for each facet and, for a direct frame
% r = (s, a, f), Sigma_r = (rho_f, rho_s, rho_a) with product 1; page 48:
% indirect frames, canonical and anticanonical Schwartz systems, when
% Sigma_r1 = Sigma_r2. Pages 49-51: Proposition for the tetrahedron (a)-f):
% 72 + 24 = 96 systems) with proof and Corollary (types (3,3,2) and (3,3,3)).
% Pages 52-56: the same for S4 via the oriented octahedron (types (4,3,2)
% and (4,4,3)), with proof of d). Pages 58-60, a new run on regular
% hyperpolygons: case n = 1 in characteristic 0 with G_0 finite, two strict
% pseudo-reflections tau_0, tau_1 in V with invariants mu_0, mu_1, beta;
% I) an invariant divisor of degree 2 forces tau_0, tau_1 of order 2, the
% ordinary regular polygons; II) otherwise the image of G_0 in Aut(C) is
% A4, S4 or A5, whence 2 <= nu_0, nu_1 <= 5. The page breaks off.
%
% Notation fixed once for the batch: his double-struck A (the ambient group
% scheme) and H as \mathbb{A}, \mathbb{H} on p. 44, plain H elsewhere as he
% writes it; the tall S and A of the symmetric and alternating groups as
% \mathfrak{S}, \mathfrak{A}; \widetilde{H} for Aut(Pi); rho_s, rho_a, rho_f
% for vertex, edge, face; R, R^+ for (direct) frames (\mathcal{R} on p. 49
% where he writes a script R); \prec for « s adjacent to f ». Underlined
% words are \emph. On p. 60 the letter for the group generated by tau_i is a
% looped 3, set Z_i with a \note.
%
% Hardest pages: 44 (long passages of fast prose, many insertions), 47 and
% 49 (successive cancellations and interlinear rewriting, slanted margins),
% 56 (a boxed cancelled block). Readings to check first: « splittée » (p. 44),
% the factor « 3 » before Aut^+(Pi) on p. 48, the type « (3, 4, 2) » added
% on p. 52, the bracket on p. 42's first line, and the signs of the
% eigenvalue equations on p. 59.

\documentclass[11pt,a4paper]{article}
\input{../preamble/grothendieck}

\folder{85}
\batch{3}
\pages{41}{60}
\dating{s.d. — le groupe « Géométrie et topologie combinatoire » (69 à 102) est daté 1976-[vers 1986]}
\watermark{Édition de démonstration}
\foldertitle{2-polyèdres et 1-polygônes : notes manuscrites (s.d.)}

\begin{document}

\section{Systèmes de Schwartz dans $\mathfrak{A}_4$, $\mathfrak{S}_4$, $\mathfrak{A}_5$, et classification des hyperpolygônes réguliers en car.\ zéro}
\note{Titre de sa main sur la couverture (page 41), avec au crayon
«~(15 p)~» : les pages 42 à 56.}

\page{42}
\[
\begin{array}{rlll}
4 & \mathfrak{A}_4 & 3 & (3, 3', 2)\,(3', 2, 3)\,(2, 3, 3')\
  \bigl[(2, 3', 3)\,(3', 3, 2)\,(3, 2, 3')\bigr] \\
  & & 1 & (3, 3, 3)\ \bigl[(3', 3', 3')\bigr] \\[4pt]
9 & \mathfrak{S}_4 & 6 & (4, 3, 2)\,(3, 2, 4)\,(2, 4, 3)\,(2, 3, 4)\,(3, 4, 2)\,(4, 2, 3)
  \ \text{standard} \\
  & & 3 & (4, 4, 3)\,(4, 3, 4)\,(3, 4, 4) \\[4pt]
32 & \mathfrak{A}_5 & 6 & (5, 3, 2)\,(3, 2, 5)\,(2, 5, 3)\,(2, 3, 5)\,(3, 5, 2)\,(5, 2, 3)
  \ \text{standard } \alpha \\
  & & 6 & (5', 3, 2)\,(3, 2, 5')\,(2, 5', 3)\,(2, 3, 5')\,(3, 5', 2)\,(5', 2, 3)
  \ \text{id. } \alpha' \\
  & & 3 & (5, 5, 3)\,(5, 3, 5)\,(3, 5, 5) \\
  & & 3 & (5', 5', 3)\,(5', 3, 5')\,(3, 5', 5') \\
  & & 6 & (5, 5', 3)\,(5', 3, 5)\,(3, 5, 5')\,(3, 5', 5)\,(5', 5, 3)\,(5, 3, 5') \\
  & & 6 & (5, 5', 2)\,(5', 2, 5)\,(2, 5, 5')\,(2, 5', 5)\,(5', 5, 2)\,(5, 2, 5') \\
  & & 1 & (5, 5, 5) \\
  & & 1 & (5', 5', 5') \\[4pt]
\hline
45 & & &
\end{array}
\]
\marginal{(ou 16 si on regarde modulo action de $\mathfrak{S}_5$)}
\note{Tableau des types numériques de systèmes de Schwartz. Les effectifs
de gauche (4, 9, 32, total 45) sont dans la marge ; ceux du milieu sont
écrits contre un trait vertical épais, et les paires de lignes de
$\mathfrak{A}_5$ sont réunies par des accolades. Dans le crochet de la
première ligne, un «~standard~» est biffé et le premier triplet est
surchargé ; sa lecture $(2, 3', 3)$ est incertaine. Dans «~id.~$\alpha'$~»,
«~id.~» est une lecture probable.}

\emph{Question}~: un syst.\ de Schwartz d'un groupe fini est-il connu
\add{mod conjugaison} (quand \uncertain{on} connaît chacun de ses termes
\uncertain{mod} conjugaison~? OK dans $\mathfrak{A}_4$, $\mathfrak{S}_4$,
$\mathfrak{A}_5$ --- regarder groupes diédraux, et $\mathfrak{S}_5$\ldots

\page{43}
Sous-groupes finis de \struck{\ill{}}
$\underline{\mathrm{Aut}}_X(\ldots) = \mathbb{A}$ ($X$ droite projective
sans composantes \add{de car.\ 0})

1) Sous-groupes \struck{$\Gamma^{+}$}\,$\Gamma^{+}, \Gamma \subset \mathbb{A}$
qui sont des formes de $\mathfrak{A}_4$ ou $\mathfrak{S}_4$. Se
correspondent biunivoquement
\[
\Gamma^{+} \longmapsto \Gamma = \mathrm{Norm}_{\mathbb{A}}(\Gamma^{+})
\qquad \Gamma^{+} = [\Gamma, \Gamma]
\]
\emph{NB} $\mathrm{Aut}(\mathfrak{A}_4) \simeq \mathrm{Aut}(\mathfrak{S}_4)
\simeq \mathfrak{S}_4$. On connaît la cat.\ des ensembles~: \ill{} \ill{},
celle des formes de $\mathfrak{A}_4$, celle des formes de $\mathfrak{S}_4$,
\struck{\ill{}} équivalences

\begin{tikzcd}[column sep=small, row sep=small]
F(\mathfrak{S}_4) \arrow[rr, "\approx"] \arrow[dr, no head, "\wr"'] & &
F(\mathfrak{A}_4) \arrow[dl, no head, "\wr"] \\
& (\mathrm{Ens})_4 &
\end{tikzcd}

\note{À droite de $F(\mathfrak{A}_4)$, un premier «~$(\mathrm{Ens})_4$~» est
biffé ; il est récrit sous les deux, relié par deux traits portant le signe
d'isomorphisme.}
\[
\begin{array}{l}
E \in (\mathrm{Ens})_4 \longmapsto \Gamma = \mathfrak{S}_E,\
\Gamma^{+} = \mathfrak{S}_E^{+} \\
\Gamma^{+} \in F(\mathfrak{A}_4) \longmapsto \Gamma = \mathrm{Aut}(\Gamma^{+}),\
E = \text{ens.\ des ss-groupes d'indice 4 (d'ordre 3) de } \Gamma^{+} \\
\Gamma \in F(\mathfrak{S}_4) \longmapsto \Gamma^{+} = [\Gamma, \Gamma],\
E = \text{ens.\ des ss-groupes d'ordre 6 (d'indice 4) de } \Gamma
\end{array}
\]
\note{Les deux dernières descriptions de $E$ ont une suite en retrait à
droite : «~ou $\Gamma \simeq$ ens.\ des ss-groupes $\gamma^{+}$ d'ordre 3
de $\Gamma$ [$\gamma^{+} = [\gamma, \gamma]$,
$\gamma = \mathrm{Norm}_{\Gamma}(\gamma^{+})$]~».}

Donc sur un schéma $S$
\[
\mathrm{Rev}_4(S) \simeq F(S, \mathfrak{S}_4) \simeq F(S, \mathfrak{A}_4)
\simeq \mathrm{Tors}(\mathfrak{S}_{4S}) \simeq
\]
\emph{cubes orientés sur $S$}

\emph{Re NB} \uncertain{Si on est} \uncertain{en car.} 0 (\uncertain{inutile}
?) si $\Gamma^{+}$ ou $\Gamma \subset$ \struck{$X$, alors}
\add{$G = \underline{\mathrm{Aut}}\,X$}, alors [tout automorphisme de
$\mathbb{A}$ est int., $\mathbb{A} \xrightarrow{\sim}
\mathrm{Aut}(\mathbb{A})$, et] les automorphismes de $(\mathbb{A}, \Gamma^{+})$
ou $(\mathbb{A}, \Gamma)$ qui normalisent $\Gamma^{+}$ ou $\Gamma$ proviennent
de $\Gamma$~: $\Gamma \xrightarrow{\sim} \underline{\mathrm{Aut}}_{(X,
\Gamma)}$, $\Gamma \xrightarrow{\sim} \underline{\mathrm{Aut}}_{(X,
\Gamma^{+})}$. Donc la cat.\ des couples $(\Gamma, X)$ $(\Gamma^{+}, X)$ est
équivalente à celle des torseurs sous $\mathfrak{S}_{4S}$, i.e.\ aux cat.\
précédentes --- donc aussi à celle des \emph{cubes orientés} sur $S$, ou
celle des \emph{tétraèdres} (pas \struck{\ill{}} orientés) sur $S$.
Dans ce
\note{Au-dessus de $\xrightarrow{\sim}$ dans le crochet, un petit mot,
peut-être «~int.~». La lettre suivant $\underline{\mathrm{Aut}}$ à la
première ligne de la page est surchargée et se lit mal.}

\page{44}
dictionnaire, $\Gamma$ s'interprète comme le \add{\ill{}} Groupe des
automorphismes du tétraèdre, ou des cubes orientés, $\Gamma^{+}$ comme le
groupe des autom.\ \emph{orientés} du tétraèdre, $\Gamma$ \uncertain{comme}
celui des automorphismes \ill{} des «~bi-orientés~» des cubes)~\ldots

2. On s'intéresse \struck{\ill{}} aux systèmes
$(\struck{H \subset G,\ \ill{}}\ \add{\mathbb{A}, \mathbb{H},}\ \rho, \rho', \rho'')$,
où $\rho, \rho', \rho''$ sont des sections \struck{\ill{}} de $G$, telles que
$\rho\rho'\rho'' = 1$, qui engendrent un des sous-groupes finis
\add{$\mathbb{H}$} d'un \ill{} \ill{} \ill{} (i.e.\ $\simeq \mathfrak{A}_4$,
$\mathfrak{S}_4$ ou $\mathfrak{A}_5$). Donc $\mathbb{H}$ sera une forme
\add{de $G$} \uncertain{splittée} de l'un de ces groupes (mais $X$ n'est pas
\uncertain{splitté} pour autant\ldots). Un \struck{\ill{}} \ill{} ainsi
(\ill{} \ill{} \ill{}) que la donnée $\rho, \rho', \rho''$ correspond à un
tétraèdre \add{comb.} \struck{\ill{}}, un cube orienté \add{comb.} (\ill{} un
\uncertain{icosaèdre} comb., cf plus bas) qui définit $X$ (\ill{} \ill{} de
\ill{} \uncertain{isosaèdre}\ldots), qu'on \ill{}. Et on se propose de trouver
les \uncertain{syst.} de trois gén.\ $\rho, \rho', \rho''$ du groupe
\struck{\ill{}} «~abstrait~» $\mathbb{H}$ des automorphismes orientés des
\struck{\ill{}} tétraèdres, ou des cubes\ldots\ \struck{\ill{}}
\marginal{\emph{NB} $\mathbb{H}$ résulte de $\rho, \rho', \rho''$, et $\rho''$
résulte de $\rho, \rho'$ (et $\rho$ de $\rho', \rho''$, $\rho'$ de
$\rho'', \rho$\ldots)}

Notons qu'un automorphisme de $(G, \rho, \rho', \rho'')$ est un autom.\ de
$(G, \mathbb{H})$ qui commute à $\rho, \rho', \rho''$ i.e.\ à $\mathbb{H}$, donc
\struck{est} trivial (car \uncertain{le centr.} de $\mathbb{H}$ se réduit au
groupe unité). Donc \ill{} \ill{} catégorie \emph{rigide}~: ses objets n'ont
pas d'autom.\ \ill{}, i.e.\ sont \emph{uniques} à isom.\ unique près, pour les
classes d'isom.\ données.
\note{Les lettres $\mathbb{A}$, $\mathbb{H}$ sont écrites en capitales doubles,
d'une encre plus appuyée, sur une première rédaction biffée «~$(H \subset G,
\ldots$~» ; le $\mathbb{H}$ ajouté au-dessus de «~sous-groupes finis~» est du
même trait.}

\page{45}
\emph{Rmq} Comme premier invariant d'un système $\rho, \rho', \rho''$, il y a
les \add{triplets des} ordres $(\nu, \nu', \nu'')$ de $\rho, \rho', \rho''$.
\struck{Les} Les op.\ de $\mathfrak{S}_3$ sur $\mathbb{H}^3$ correspondent
évidemment aux op.\ de $\mathfrak{S}_3$ sur $\mathbf{N}^3$. Notons
$\nu, \nu', \nu'' \geq 2$ et
\[
\begin{array}{lll}
\mathrm{I} & 2 \leq \nu, \nu', \nu'' \leq 3 & \text{cas } \mathfrak{A}_4 \\
\mathrm{II} & 2 \leq \nu, \nu', \nu'' \leq 4 & \text{cas } \mathfrak{S}_4 \\
\mathrm{III} & 2 \leq \nu, \nu', \nu'' \leq 5 & \text{cas } \mathfrak{A}_5
\end{array}
\]
\note{Les trois lignes sont réunies par une accolade.}

On verra que dans le premier cas, le triplet \struck{\ill{}} contient 3, dans
le deuxième il contient 4, dans le troisième il contient 5 --- donc le
triplet \ill{} \ill{} indique déjà dans lequel des cas I, II, III on se
trouve (en \uncertain{regroupant} \ill{} \ill{} \ill{} 3, 4 ou 5). Dans les
cas I, II, on voit que $(\rho, \rho', \rho'')$ est déterminé \ill{} \ill{} par
sa condition numérique. Dans le cas icosaèdr.\ III, il n'en est plus ainsi~:
fait ainsi (\ill{} \ill{} par l'équivalence modulo $\mathfrak{S}_5$ qui
\ill{}) --- il faut préciser un \struck{\ill{}} la condition numérique [en
distinguant les deux classes de conjugaison d'él.\ d'ordre 5 dans
$\mathfrak{A}_5$\ldots]

3) Cas du $\mathfrak{A}_4$.

\emph{Théorème} Soit $\Pi$ un tétraèdre \add{comb.} orienté, $H$ le groupe de
ses automorphismes \add{(orientés)}. À tout sommet $s \in S$ de $\Pi$ est
associé un $\rho_s \in H$, él.\ d'ordre 3, et on trouve ainsi une
\uncertain{injection} de $S$ \struck{\ill{}} l'ens.\ des \add{8} éléments
d'ordre 3 de $H$, soit $H_3^{+}$ son image, de sorte
\note{À partir de «~Théorème~», la fin de la page et le haut de la page 47
jusqu'à la parenthèse fermant a) sont barrés de longs traits obliques.}


\page{46}
Notons que si $H$ est un groupe abstrait, $\mathrm{Tr}(H)$ l'ens.\ des
$(\rho, \rho', \rho'') \in H^3$ tels que $\rho\rho'\rho'' = 1$, le groupe
$\mathfrak{S}_3$ opère dedans de façon évidente, l'orbite de
$(\rho, \rho', \rho'')$ \struck{$(\rho, \rho', \rho'')$ \ill{}
$(\rho', \rho'', \rho)$} \ill{} formée des
\[
(\rho, \rho', \rho'')\,(\rho', \rho'', \rho)\,(\rho'', \rho, \rho')\,
(\rho''^{-1}, \rho'^{-1}, \rho^{-1})\,(\rho'^{-1}, \rho^{-1}, \rho''^{-1})\,
(\rho^{-1}, \rho''^{-1}, \rho'^{-1})
\]
De cette façon, le groupe $\mathfrak{S}_3$ opère sur la cat.\ des
«~groupes \uncertain{pythagoriciens} avec présentation de Schwartz~» ou
simplement systèmes de Schwartz.

On est intéressé à tenir compte de cette opération, et notamment à
déterminer quels sont les transformés \add{(par op.\ de $\mathfrak{S}_3$)}
d'un $(H, \rho, \rho', \rho'')$ qui lui sont isomorphes.
\note{«~Schwartz~» est son orthographe, ici et plus loin.}

\emph{Rmq} La notion d'isom.\ qui nous intéresse est celle définie par le
sch.\ en groupes ambiant $\mathbb{A}$, ce qui s'explicite en termes
élémentaires ainsi~:
\begin{itemize}
\item[a)] Cas $\mathfrak{A}_4$~: isomorphie par op.\ de $\mathfrak{S}_4$
\add{\ill{} $\mathfrak{A}_4$} (\emph{pas} $\mathfrak{A}_4$). On peut
évidemment se proposer de faire aussi la classification plus \emph{fine}
\ill{} op.\ de $\mathfrak{A}_4$.
\item[b)] Cas $\mathfrak{S}_4$~: isom.\ par autom.\ int.
\item[c)] Cas $\mathfrak{A}_5$~: isom.\ par autom.\ int. Il y a une
classification moins fine, modulo autom.\ de $\mathfrak{S}_5$ sur
$\mathfrak{A}_5$, qui donne unicité «~moins~» de classe d'isomorphie.
\end{itemize}

\page{47}
que $H_3^{-} = H_3 \setminus H_3^{+}$ est l'ens.\ des
$(\rho_s^{-1})_{s \in S}$, ou encore l'ens.\ des $\rho_f$ ($f \in F$ ens.\ des
faces) \struck{Notons que si $s, f$ \ill{} \ill{}}
\marginal{\emph{NB} $s \mapsto f$ \ill{} \ill{} $\rho_f = \rho_s^{-1}$}

a) Pour toute \add{arête} \struck{\ill{}} $\{s, t\}$ de $\Pi$, le triple
\[
(\rho_t^{-1}, \rho_s, \rho_t^{-1}\rho_s)
\]
engendre $H$, et est de type numérique $(3, 3, 2)$, \add{les trois}
\struck{Les} él.\ d'ordre 2 de $H$ étant identifiés aux \add{paires}
\struck{\ill{}} d'arêtes «~opposées~» de $\Pi$ (comme \add{\ill{}}
\struck{\ill{}} \ill{} à \ill{} joignant leurs milieux).

3. De façon générale, si $\Pi$ est un polyèdre \add{régulier,
\uncertain{orientable}} comb.\ \struck{orienté}, et
$H = \mathrm{Aut}^{+}(\Pi)$, \struck{\ill{}} \add{associé}, \ill{} facette
$\varphi \in \mathrm{Fac}(\Pi) = S \amalg A \amalg F$, un él.\
$\rho_\varphi \in H$ \add{(en utilisant l'orientation)}, de sorte que
\ill{} \add{\emph{direct}} repère $r = (s, a, f)$, le triple
\[
\Sigma_r = (\rho_f, \rho_s, \rho_a) \in H^3
\]
satisfasse $\rho_f \rho_s \rho_a = 1$ (\emph{NB} $\rho_f\rho_s = \rho_a$ et
$\rho_a^2 = 1$) ou encore $\rho_s \rho_a \rho_f = 1$. On note que tous ces
systèmes de générateurs sont conjugués \uncertain{sous}
$\mathrm{Aut}(\Pi)$\struck{, mais pas nécessairement sous
$H = \mathrm{Aut}^{+}(\Pi)$} \struck{--- il faut distinguer les $\Sigma_r$
suivant que $r$ est \ill{} à l'une ou l'autre orientation de $\Pi$.
\struck{\ill{}} La question est donc quand il est vrai que $\Sigma_r$ et
$\Sigma_{r'}$ \uncertain{conjugués} sous $H$} \add{\ill{}}
\struck{\ill{} $r, r'$ \ill{} orientation \ill{}, p.ex.}
\uncertain{prenant} l'orientation \ill{}, \ill{} \ill{}
$\rho_f^{-1}\rho_s^{-1}\rho_a^{-1} = \rho_a$ \ill{} \ill{} au lieu de
$\rho_f \rho_s \rho_a$, on trouverait
\marginal{\emph{NB} $\Sigma_r$ ne dépend pas du choix préalable d'une
orientation, car $r$ en \uncertain{définit} une}
\note{Deux passages de la seconde moitié de la page sont biffés, l'un par des
traits obliques, l'autre encadré puis barré ; la lecture de leurs limites est
incertaine.}

\page{48}
d'autres systèmes de gén.
\[
\Sigma_{r'} = (\rho_f^{-1}, \rho_s^{-1}, \rho_a = \rho_a^{-1})
\]
associés aux repères $r' = (s, a, f)$ \emph{indirects}. Les $\Sigma_{r'}$ sont
conjugués entre eux sous $H$, et les $\Sigma_r$ et $\Sigma_{r'}$ sont
conjugués sous $\mathrm{Aut}\,\Pi$, mais pas nécessairement sous
$H = \mathrm{Aut}^{+}(\Pi)$ --- par exemple pour qu'il en soit ainsi, il faut
que les $\rho_s$ et $\rho_s^{-1}$ (ou $\rho_f$ et $\rho_f^{-1}$) soient
\ill{} conjugués (cf icosaèdre). \emph{Systèmes de Schwartz canoniques} et
\emph{anticanoniques} associés au polyèdre régulier orienté.

Quand notamment
\[
\Sigma_{r_1} = \Sigma_{r_2}
\]
($r_1$ et $r_2$ \struck{directs} repères donnés) --- disons \ill{} \ill{}
orientation, i.e.
\[
\rho_{s_1} = \rho_{s_2},\quad \rho_{f_1} = \rho_{f_2},\quad
\rho_{a_1} = \rho_{a_2}.
\]
\struck{Il semble} Posons $r_2 = g \cdot r_1$ ($g \in \mathrm{Aut}\,\Pi$ bien
déterminé), alors
\[
\Sigma_{r_2} = g \cdot \Sigma_{r_1}
= (\mathrm{int}(g)\rho, \mathrm{int}(g)\rho', \mathrm{int}(g)\rho'')
\]
si $\Sigma_{r_1} = (\rho, \rho', \rho'')$. Donc $\Sigma_{r_2} = \Sigma_{r_1}$
ssi $g$ \add{commute} au \struck{\ill{}} \uncertain{centre} de
$\mathrm{Aut}^{+}\Pi$ \add{dans $\mathrm{Aut}\,\Pi$}. Dans le cas où
$\mathrm{Aut}\,\Pi = \uncertain{3} \cdot \mathrm{Aut}^{+}(\Pi)$ \add{(cas des II
et III)}, \ill{} \uncertain{veut} donc que \struck{\ill{}} tous les
$\Sigma_r$ \uncertain{conjugués} \ill{} $\mathrm{Aut}^{+}(\Pi)$ (pour les $r$
directs et inverses)
\note{Le facteur devant $\mathrm{Aut}^{+}(\Pi)$ est lu «~3~», ce qui ne
s'accorde pas avec un sous-groupe d'indice 2 ; peut-être un $\mathbf{Z}/2$ ou
un $\times$ mal formé. Transcrit tel que lu.}

\page{49}
\emph{Proposition} Soit $\Pi$ un tétraèdre orienté, $\mathcal{R}$
\add{$(\mathcal{R}^{+})$} l'ens.\ de ses repères \add{rep.}
(\struck{directs}), $S, A, F$ l'ens.\ des sommets, \struck{\ill{}} arêtes,
faces --- d'où $\rho_s, \rho_a, \rho_f$ si $s \in S$, $a \in A$, $f \in F$.
Soit $\Phi = S \amalg A \amalg F$ l'ens.\ des facettes.
\struck{a) L'application $\ill{}\ 4\ \ill{}\ \rho_s$}
\begin{itemize}
\item[a)] L'application $\varphi \mapsto \rho_\varphi$ de $S \amalg F$ dans
$H = \mathrm{Aut}^{+}\Pi$ est injective, et son image est l'ens.\ des 8~él.\
d'ordre 3 de $H$. De plus, si $s \in S$, $f \in F$ sont
\uncertain{opposés}~: $\rho_s\rho_f = \rho_f\rho_s = 1$.
\item[b)] L'application $a \mapsto \rho_a$ de $A$ dans $H$ induit une
application de degré 2 de $A$ dans l'ens.\ \add{(\uncertain{\ill{}})}
${}_{3}H$ des él.\ d'ordre 3 de $H$ [\emph{NB} deux arêtes, si
\uncertain{elles ont} même image, \ill{} \ill{} opposées] \uncertain{ainsi}
que \ill{} \ill{} des parties \uncertain{disjointes} de $S$, \ill{} de $F$
\item[c)] L'application $r \mapsto \Sigma_r$ de $\mathcal{R}$ dans l'ens.\ des
\struck{\ill{}} syst.\ de Schwartz de $H$ est \uncertain{surjective}
[\ill{} \ill{} le commutant \ill{} $H$ dans
\add{$\widetilde{H} =$} $\mathrm{Aut}(\Pi)$ est $1$] l'image formée d'une
classe sous $\widetilde{H}$ (\struck{\ill{}} \add{\ill{}} \ill{} torseur sous
$\widetilde{H}$) et deux classes sous $H$, de type numérique $(3, 3, 2)$.
\item[d)] Si \ill{} a un syst.\ de Schwartz canonique [ou anticanonique
$(\rho_f', \rho_s', \rho_a')$]
$(\rho_f, \rho_s, \rho_a)$, les \add{6} systèmes qu'il engendre, par
\uncertain{action} de $\mathfrak{S}_3$ \ill{}, sont deux à deux
\struck{\ill{}} \ill{} \ill{} \ill{} \ill{} \ill{} \ill{}
\struck{$(\rho, \rho', \rho'')$ \ill{} \ill{}} \ill{} types numériques
\ill{}, aux \struck{$(\rho', \rho'', \rho)\,(\rho'', \rho, \rho')$}
\ill{} qui sont de types numériques comme
\[
(\rho_f, \rho_s, \rho_a),\ (\rho_s^{-1}, \rho_f^{-1}, \rho_a^{-1} = \rho_a)
\quad \text{i.e.\ } \Sigma_r \text{ et } \Sigma_{r'}
\]
\struck{sont conjugués} \add{\ill{}} sous \add{$\rho_f'$, $\rho_s'$,
$\rho_a'$} conjugués \add{\uncertain{\ill{}} par} $H$, \struck{mais conjugués
par} $\widetilde{H}$. \struck{(car $r, r'$ définissent des orientations
opposées) \ill{} donc \ill{} conjugaison \ill{} \ill{}}
($r'$ \uncertain{repère opposé} de $r$, \struck{il \ill{}} \add{\ill{}
orientation!})
\end{itemize}
\marginal{$24 = 12 + 12$ syst.\ de Schwartz}
\marginal{[cela fait 72 $= 36 + 36$ \ill{} \ill{} sous $\mathfrak{S}_3$,
\ill{} \ill{} 36 syst.\ de Schwartz \struck{\ill{}} \ill{} \ill{}
$\mathfrak{S}_3 \times H$\ldots]}
\note{La fin de la page, à partir de «~sous $\rho_f'$~», est très surchargée :
biffures successives, ajouts interlinéaires et, en bas à droite, un petit
croquis de deux triangles reliés par des flèches, non reproduit. La note
marginale de gauche est écrite en biais et se lit mal ; seuls «~72~»,
«~36~», «~$\mathfrak{S}_3$~» et «~$H$~» sont sûrs.}

\page{50}
\begin{itemize}
\item[e)] Soient \struck{\ill{} \ill{}} $s, s', s'' \in S$ formant une face
$f$, et induisant \struck{\ill{}} \add{par leur} ordre \struck{opposé} une
orientation de $f$ \uncertain{opposée} de celle \uncertain{induite} par
$\Pi$ \struck{$a_s, a_{s'}, a_{s''}$} (i.e.\ si $\ill{}$ est le sommet
opposé à $f$, $\rho_f$ \ill{}
\struck{$s' = \rho_f s$, $s'' = \rho_f s'$} $s \mapsto s' \mapsto s''
\mapsto s$) alors
\[
(\rho_s, \rho_{s'}, \rho_{s''})
\]
est un système de Schwartz de type $(3, 3, 3)$, et si $f, f', f''$ sont les
faces \ill{} \ill{}, alors
\[
(\rho_{f''}, \rho_{f'}, \rho_f)
\]
est un syst.\ de Schwartz de type $(3, 3, 3)$.
\struck{Les 24} \add{Chacun des \uncertain{deux} paquets de 12} syst.\ de
Schwartz $\{\rho_s, \rho_{s'}, \rho_{s''}\}$,
$\{\rho_f, \rho_{f'}, \rho_{f''}\}$ forment \struck{deux} une orbite sous
$H$ (donc un torseur sous $H$) et leur réunion est une \ill{} sous
$\widetilde{H}$ (donc un torseur sous $\widetilde{H}$).
\marginal{\struck{12} syst.\ de Schwartz}
\marginal{\emph{NB} \add{Dans un} \ill{} sous-groupe de $\ill{} \simeq
\mathfrak{A}_4$ \ill{} \ill{} $H$ \ill{} \ill{} d'ordre 3 \struck{\ill{}}
\ill{} \ill{}}
\marginal{12 syst.\ de Schwartz}
\marginal{\emph{NB} les \uncertain{paquets} de 12 \ill{} \ill{} \ill{} les
\ill{} \ill{} $\mathfrak{S}_3$, \ill{} \ill{} de $\mathfrak{S}_3$ \ill{}
les deux paquets}
\item[f)] Les \struck{72} \add{$72 + 24 = 96$} systèmes de Schwartz
précédents \uncertain{épuisent} \struck{\ill{}} $H$ \ill{} les seuls.
\end{itemize}
\note{Deux traits horizontaux séparent ici l'énoncé de sa démonstration.
Les notes marginales de gauche, écrites en biais et en partie biffées, ne se
lisent que par fragments.}

\emph{Prouvons} f). \uncertain{Tenu} d'abord, comme les él.\ $\neq 1$ de $H$
sont d'ordre 2 et 3, et que les éléments d'ordre 2 sont contenus dans un
\ill{} sous-groupe \ill{} \add{$V$} d'indice 2, on voit que un syst.\ de
Schwartz ne peut contenir deux él.\ d'ordre 2, donc il contient au moins
2~él.\ d'ordre 3. \struck{Soit $(\rho, \rho', \rho'')$}

\page{51}
À \struck{\ill{}} op.\ près de $\mathfrak{S}_3$, \emph{NB} donc que le
système est de la forme $(\rho, \rho', \rho'')$, avec $\rho, \rho'$ d'ordre
3, donc chacun de la forme $\rho_s$ ou $\rho_f$
\begin{itemize}
\item[1)] $\rho = \rho_f$, $\rho' = \rho_s$~: il faut $s \prec f$ (car si
$s$ opposé à $f$, $\rho_s = \rho_f^{-1}$ donc $\rho, \rho'$
n'engendrent pas $H$), donc on est dans le cas canonique
\item[2)] $\rho = \rho_s = \rho'_{f'}$, $\rho' = \rho_f = \rho'_{s'}$~: il
faut $s \prec f$ ou $s' \prec f'$, donc on est dans le cas anticanonique
\item[3)] $\rho = \rho_s$, $\rho' = \rho_t$ \add{$s, t \in S$ ($s \neq t$)}~:
on est dans le \add{1\textsuperscript{er}} (cas de e))
\item[4)] $\rho = \rho_f = \rho'_{s'}$, $\rho' = \rho_g = \rho'_{t'}$
\add{$f, g \in F$ ($f \neq g$)}~: on est dans le 1\textsuperscript{er} cas
de e)
\end{itemize}
\marginal{12, 12, 12, 12}
\note{Un «~12~» est porté dans la marge gauche en face de chacun des quatre
cas. Dans 3) et 4), la lettre finale ($\rho_t$, $\rho_g$) est écrite sur une
autre, raturée.}

\emph{Corollaire} Les types numériques possibles des syst.\ de Schwartz pour
$H \simeq \mathfrak{A}_4$ sont
\begin{itemize}
\item[a)] $(3, 3, 2)\ (3, 2, 3)\ (2, 3, 3)$ [cas canonique et anticanonique]
\item[b)] $(3, 3, 3)$
\end{itemize}
\struck{Et} La classe de conjugaison sous
$\widetilde{H} = \mathrm{Aut}_{\mathrm{gr}}(H)$ ($\simeq \mathfrak{S}_4$) est
déterminée par le type numérique.

Donc un $H$ ($\simeq \mathfrak{A}_4$) avec syst.\ de Schwartz est de l'un des
types suivants (où $H = \mathrm{Aut}^{+}(\Pi)$, $\Pi$ un tétraèdre comb.,
donc \ill{} unique \ill{} d'orientation privilégiée)
\marginal{\emph{NB} La donnée de $H$ équivaut à celle de $\Pi$, cf plus
haut}
\begin{itemize}
\item[a)] On a un repère $r$ (unique, déterminé) de $\Pi$, et
$\Sigma = \Sigma_r$ (cas canonique)
\item[b)] On a \add{un \ill{}} \struck{\ill{}} repère
$r = \{s, \{s, s'\}, \{s, s', s''\}\}$ \add{bien déterminé} et
$\Sigma = (\rho_s^r, \rho_{s'}^r, \rho_{s''}^r)$ \ill{} \ill{}
\end{itemize}
\note{La note marginale est entourée. La dernière ligne de la page est
surchargée d'ajouts et de barres obliques ; sa fin se perd au bord.}

\page{52}
\add{\ill{}} \struck{$r$, indiquons que l'on utilise l'orientation définie
par $r$ pour définir} \struck{\ill{} de $\Pi$, et une f.}
$\rho_s, \rho_{s'}, \rho_{s''}$ \struck{\ill{} \ill{}} ou soit (bien
déterminé) $(s, s', s'')$ \ill{} él.\ distincts de $S$, de sorte que, si on
\ill{} l'orientation de $\Pi$ \struck{telle} \add{qui induise sur}
\struck{que la f.} la face $f = \{s, s', s''\}$ l'orientation opposée de
celle définie par l'ordre précédent, on ait
\[
\Sigma = (\rho_s^{\omega}, \rho_{s'}^{\omega}, \rho_{s''}^{\omega}).
\]
[\emph{NB} la donnée de $s, s', s''$ équivaut à celle d'un repère, savoir
$r = \{s, \{s, s'\}, \{s, s', s''\}\}$].
\note{Les deux premières lignes sont barrées et un fragment en est encadré ;
un long trait oblique sépare ensuite ce qui précède de la suite.}

Cas de $\mathfrak{S}_4$

\emph{Proposition} Soit $\Pi$ un octaèdre comb.\ \add{orienté}, et
$H = \add{\mathrm{Aut}^{+}(\Pi)}$ \struck{le groupe de} ses autom.\ orientés,
$R = \mathrm{Rep}(\Pi)$, $R^{+} = \mathrm{Rep}^{+}(\Pi)$
\begin{itemize}
\item[a)] Pour tout $r \in R$, soit $\Sigma_r$ le syst.\ de Schwartz
correspondant \add{de type \uncertain{$(3, 4, 2)$}}. Alors
$\Sigma_r = \Sigma_{r'}$ ssi $r$ et $r'$ sont antipodiques, et l'ens.\ des
$\Sigma_r$ ($r \in R$) est égal à l'ens.\ des $\Sigma_r$ ($r \in R^{+}$) et en
corr.\ 1-1 avec $R^{+}$ --- donc un torseur sous $H$.
\item[b)] L'application $s \mapsto \rho_s$ définit une bijection de $S$ avec
${}_4H$ (ens.\ des él.\ d'ordre 4 de $H$)
\end{itemize}
\marginal{24 d'où $6 \times 24 = 144$ syst.\ de Schwartz (en multipliant par
$\mathfrak{S}_3$)}
\note{Le chiffre du type ajouté en a) est surchargé ; la lecture $(3, 4, 2)$
est incertaine.}

\page{53}
l'application $f \mapsto \rho_f$ définit une bijection de $F$ avec ${}_3H$,
enfin l'application $a \mapsto \rho_a$ fait de $A$ un rev.\ de degré 2 de
l'ens.\ ${}_2H \setminus {}_2H^{+} = {}_2H^{-}$ (où
$H^{+} = [H, H] \simeq \mathfrak{A}_4$) formé des él.\ d'ordre 2 qui ne sont
pas dans le s-groupe invariant d'ordre 4 de $H$~; on a
$\rho_a = \rho_{a'}$ ssi $a, a'$ sont antipodiques. Les applications
$s \mapsto \rho_s$ et $f \mapsto \rho_f$ transforment antipodisme en passage
à l'inverse $\rho \mapsto \rho^{-1}$.
\begin{itemize}
\item[c)] Pour toute arête $(s, t)$ de $\Pi$, soit $f$ l'unique face telle que
$(s, \{s, t\}, f) \in R^{+}$. \struck{avec rep.} Alors
\[
(\rho_s, \rho_t, \rho_{f'}) \qquad (f' \text{ l'antipodique de } f)
\]
est un syst.\ de Schwartz de type $(4, 4, 3)$. \add{L'ens.\ de ces}
\struck{Les} syst.\ de Schwartz forment un torseur sous $H$.

Si on regarde les \add{6} transformés de $(\rho_s, \rho_t, \rho_f)$ par
$\mathfrak{S}_3$, ils se groupent en couples de 2 ayant même condition
numérique, comme $(\rho_s, \rho_t, \rho_f)$ et
$(\rho_t^{-1}, \rho_s^{-1}, \rho_f^{-1}) = (\rho_{t'}, \rho_{s'}, \rho_{f'})$,
dont les deux sont conjugués sous $H$ (en effet,
$(t', \{s', t'\}, f') \in R^{+}$ car l'antipodisme renverse l'orientation)
\item[d)] Les syst.\ de Schwartz ci-dessus épuisent tous les cas possibles.
\end{itemize}
\marginal{24 d'où $24 \cdot 3 = 72$ syst.\ de Schwartz en multipliant par
$\mathfrak{S}_3$}
\marginal{\emph{NB} les s-gr.\ $V$ de $H = \mathfrak{S}_4$ \ill{} \ill{}
d'ordre \ill{} \ill{} de 3 et 4, \ill{} $H^{+} = \mathfrak{A}_4$, \ill{} $H$
(\ill{} $\rho_s\rho_t\rho_f \notin H^{+}$!)}
\note{En c), l'énoncé porte $\rho_{f'}$ dans la formule centrée, puis
$\rho_f$ dans la discussion des transformés ; transcrit tel quel.}

\page{54}
\emph{Corollaire} Les types numériques possibles de syst.\ de Schwartz pour
$H = \mathrm{Aut}^{+}(\Pi)$ sont les suivants
\begin{itemize}
\item[a)] $(4, 3, 2)\ (3, 2, 4)\ (2, 4, 3)\ (2, 3, 4)\ (3, 4, 2)\ (4, 2, 3)$
(canoniques)
\item[b)] $(4, 4, 3)\ (4, 3, 4)\ (3, 4, 4)$
\end{itemize}
Le type numérique détermine le syst.\ de Schwartz modulo conjugaison. De
façon précise, si $\Sigma = (\rho, \rho', \rho'')$ est un syst.\ de Schwartz,
alors \uncertain{il} est dans l'un ou l'autre des cas suivants, s'excluant
mutuellement
\begin{itemize}
\item[a)] $\exists$ repère direct $r \in R^{+}$ \add{(unique)} tel que
$\Sigma$ soit égal à $\Sigma_r$ ou un de ses transformés par
$\sigma \in \mathfrak{S}_3$ ($\sigma$ unique aussi)
\item[b)] $\exists$ \struck{\ill{} \ill{}} $r \in R^{+}$,
$r = (s, \{s, t\}, f)$, \struck{et} tel que $\Sigma$ soit égal à
$(\rho_s, \rho_t, \rho_f)$ ou à un transformé par
$\sigma \in \mathfrak{S}_3^{+}$ ($\sigma$ aussi unique).
\end{itemize}

\emph{Dém.\ de} d). Les ordres des él.\ $\neq 1$ de $H$ sont $2, 3, 4$.
Montrons d'abord que un système de Schwartz contient au moins un él.\
d'ordre 4, et un él.\ d'ordre 3.

\page{55}
$\rho = \rho_s$, $\rho' = \rho_f$. Si $s \prec f$, on est dans le cas déjà
traité dans a). Sinon, on a $s \prec f'$ \add{$\rho_{f'}^{-1}$} donc
\struck{$(\rho_{f'}, \rho_s, \rho_a)$ est un syst.\ de Schwartz canonique
(où $a \in A$ est défini de façon que $(s, a, f') \in R^{+}$) donc
$(\rho_s^{-1}, \rho_{f'}^{-1} = \rho_f, \rho_a^{-1} = \rho_a)
= (\rho^{-1}, \rho', \rho_a)$ est un syst.\ canonique. Donc \ill{}
$(\rho, \rho', \rho'') = \rho = \rho_s$, $\rho$} \add{Cela donne une}
seule autre possibilité comme position relative d'un él.\ d'ordre 4 et d'un
él.\ d'ordre 3 qui \uncertain{engendrent}, celle traitée dans c) [type
$(4, 3, 4)$] \uncertain{absurde}. cqfd.
\note{Le passage biffé est encadré puis barré de traits obliques ; l'ajout
«~Cela donne une~» est relié par un trait à la suite non biffée. Dans la
marge gauche, un petit triangle dont les sommets et un côté portent $s$,
$a$ et, à l'intérieur, $f'$. La moitié inférieure de la page est blanche.}

\page{56}
1) Il contient un él.\ d'ordre 4, i.e.\ n'est pas formé d'él.\ d'ordre 2 et
3.

Il n'est \struck{pas formé de} \add{\uncertain{contenir}} deux él.\ d'ordre 2,
sinon $H$ serait un groupe diédral, ce qui n'est pas le cas. Donc il
contiendrait deux él.\ d'ordre 3, \struck{qui engendrent $H$ \ill{} \ill{}}
\struck{\ill{}} les él.\ d'ordre 3 sont dans $H^{+}$, absurde.
\note{Entre «~deux él.\ d'ordre 3,~» et «~les él.\ d'ordre 3~», un passage
de cinq lignes est encadré et biffé ; on y lit notamment «~$\rho_{f'}$,
$\rho_{g}$~», «~$H\ (\simeq \mathfrak{S}_4)$~», «~$\Pi_{123}$, et $\Pi_{234}$
ou $\Pi_{324}$~», «~(\ldots\ conjugaison)~», et le reste ne se lit pas.}

2) Il contient aussi un él.\ d'ordre 3, i.e.\ n'est pas formé seulement
d'él.\ d'ordre 4 et 2. Car dans ce cas, comme il ne contient \uncertain{pas}
2 él.\ d'ordre 2, il contiendrait au moins un autre él.\ d'ordre 4.
\struck{Soient} \add{\emph{NB}} $\rho, \rho'$ dans $(\rho, \rho', \rho'')$
\uncertain{que} $\rho, \rho'$ sont d'ordre 4 (quitte à utiliser
$\mathfrak{S}_3$), donc de la forme $\rho_s, \rho_t$. \struck{Mais}
\add{Évidemment} $t \neq s, s'$ (antipodique de $s$), donc $(t, s)$ est une
arête, on est dans le cas traité en c), où on a un él.\ d'ordre 3.

3) Il reste \uncertain{à examiner} le cas des 2 él.\ d'ordre 4 \add{traité}
\ill{} \uncertain{regardons} le cas où \ill{} a un seul él.\ d'ordre 4, un
d'ordre 3, soit $(\rho, \rho', \rho'')$ \ill{}


\section{Hyperpolygônes réguliers}
\note{Titre de sa main sur la couverture (page 57), avec au crayon
«~(16 p)~» : les pages 58 à 73, dont seules 58 à 60 sont dans ce lot.}


\page{58}
Cas $n = 1$, \uncertain{sur} \uncertain{corps} \add{\uncertain{alg.\ clos}}
car.\ 0, $G_0$ fini
\[
\mu_0 = \lambda_0 - 1,\quad \mu_1 = \lambda_1 - 1,\quad \beta = \alpha + 2
\quad \text{invariant}
\]
\begin{itemize}
\item[1)] $\mu_0, \mu_1 \neq 0$ i.e.\ $\lambda_0, \lambda_1 \neq 1$
($\tau_0, \tau_1$ d'ordre fini),
$\lambda_0, \lambda_1$ racines de l'unité (d'ordres
$\underline{\nu}_0, \underline{\nu}_1 \geq 2$)
\item[2)] $\langle \pi, \pi' \rangle = \rho \neq 0$
\add{$= \mu_0\mu_1 - \beta$} i.e.\ $\mathcal{E} \simeq V \oplus k\pi$.
On peut se borner à $V$~:
\end{itemize}
Système de deux ps.\ réflexions strictes $\tau_0, \tau_1$ dans $V$, avec
conditions de non dégénérescence~:
\begin{itemize}
\item[a)] $\langle a_0, a_1' \rangle \neq 0$
\item[b)] $a_0, a_1$ lin.\ indép.
\end{itemize}
\note{Le signe devant $\geq 2$ est écrit sur un autre, raturé.}

I) La représentation dans $C = \check{\mathbf{P}}(V)$ admet un \uncertain{div.}
de degré 2 invariant

Il ne peut être réduit à un pt, car s'il y a une droite \add{$\subset V$}
\emph{invariante} par $G_0$, il y a un supplém., donc $G_0$ est commutatif
$\subset k^{\times} \times k^{\times}$, absurde.

Donc $V = D \oplus D'$, et pour le $=$ \uncertain{raison} que dessus, $D, D'$
ne sont pas invariants par $\tau_0, \tau_1$~; il faut \struck{a)
$\tau_0$ invariant $D, D'$, $\tau_1$ \ill{}} \add{\uncertain{si} $\tau_i$
invariant $D$, i.e.\ que ou bien $a_i'$ est nul sur} $D$, ou bien
$a_i \in D$. Dans le premier cas, on aurait $a_i \in D'$, dans le deuxième
$a_i'$ nul sur $D'$. Ainsi, OPS $a_i \in D$, $a_i'$ nul sur $D'$, donc
\[
D = k a_i, \qquad D' = \mathrm{Ker}\, a_i'
\]
Considérons alors \add{$j \neq i$, $j \in \{0, 1\}$} \struck{$j \neq i$
\ill{}}, donc $\tau_j D \subset D'$, $\tau_j D' \subset D$,
\[
a_i + \underbrace{\langle a_i, a_j' \rangle}_{\struck{\ill{}}\ \beta} a_j
\in D' \quad \text{i.e.} \quad
\underbrace{\langle a_i, a_i' \rangle}_{\mu_i}
+ \underbrace{\langle a_i, a_j' \rangle \langle a_j, a_i' \rangle}_{\beta} = 0
\]
donc $\beta \neq 0$, donc quitte à échanger $\tau_0, \tau_1$, OPS que $i = 0$,
$j = 1$, $\beta = -\mu_0$. Exprimons aussi $\tau_1 D' \subset D$ i.e.\
$\tau_1^2 D \subset D$, $\tau_1^2 = \mathrm{id} + (2 + \mu_1)\, a_1' \otimes a_1$
\note{Dans le premier accolade, sous $\langle a_i, a_j' \rangle$, un mot biffé
puis «~$\beta$~» ; la lecture de l'accolade est incertaine.}

\page{59}
on trouve $a_0 + (2 + \mu_1)\langle a_0, a_1' \rangle a_1 \in D'$ i.e.
\[
\underbrace{\langle a_0, a_0' \rangle}_{\mu_0}
+ (2 + \mu_1)\underbrace{\langle a_0, a_1' \rangle \langle a_1, a_0' \rangle}_{\beta}
= 0
\]
i.e.\ $\mu_0 + (2 + \mu_1)\beta = 0$ i.e.\ \struck{$(2 + \mu_1)\beta$ (comme
$\beta$} (comme $\mu_0 = -\beta$) $(1 + \mu_1)\beta = 0$ i.e.\ (comme
$\beta \neq 0$), $1 + \mu_1 = 0$ i.e.\ \struck{$\mu_1 = -1$ i.e.\
$\lambda_1 - 1 =$} $\lambda_1 = 0$, absurde.

Ainsi $\tau_0, \tau_1$ échangent $D_0, D_1$. Prouvons \struck{\ill{}}
\add{que} $\mu_0 = \mu_1 = -2$, i.e.\ $\lambda_0 = \lambda_1 = -1$, i.e.\
$\tau_0, \tau_1$ \emph{d'ordre 2} --- i.e.\ qu'on est dans le cas des
polygônes réguliers ordinaires. \struck{Pour ce qui précède, \ill{} \ill{}}
\add{Prenons coordonnées adaptées à} $D, D'$
\struck{Notons $D = k a_0$, $D' = k a_1$ donc}
$e_0 \in D$, $e_1 = \tau_0 e_0$, on aura, si $\tau_0 e_1 = \xi e_0$,
$\tau_1 e_0 = \eta e_1$, \add{$\tau_1 e_1 = \zeta e_0$}
\[
\tau_0 = \begin{pmatrix} 0 & \xi \\ 1 & 0 \end{pmatrix}, \qquad
\tau_1 = \begin{pmatrix} 0 & \zeta \\ \eta & 0 \end{pmatrix}
\]
équations aux valeurs propres
\[
\lambda^2 + \xi = 0, \qquad \lambda^2 + \zeta\eta = 0
\]
et comme cette équation a la \struck{\ill{}} \add{\uncertain{valeur propre}}
1, on doit avoir l'autre valeur propre $\lambda_0$ (resp.\ $\lambda_1$) égale
à $-1$ OK.
\note{Les signes dans les équations aux valeurs propres sont transcrits tels
qu'écrits : avec $\lambda^2 + \xi = 0$ la valeur propre 1 ne s'obtient pas
sans $\xi = -1$, que la page ne dit pas. Les lettres $\xi$ et $\zeta$ sont
proches dans sa main.}

II) Il n'y a pas dans $\check{\mathbf{P}}(V)$ de div.\ de degré 2 invariant.

Alors on sait que l'image de $G_0$ dans $\mathrm{Aut}(C) \simeq
\mathrm{GP}(1, k)$ est un des trois groupes $\mathfrak{A}_4$,
$\mathfrak{S}_4$, $\mathfrak{A}_5$. Considérons leurs images inverses dans
$\mathrm{Sl}(V)$, $\widetilde{\mathfrak{A}}_4$,
$\widetilde{\mathfrak{S}}_4$, $\widetilde{\mathfrak{A}}_5$ \add{\ill{}
$G_0$}. On a donc
\[
G_0 \subset \widetilde{G}_0 \cdot k^{\times}
\]
\note{Sous $k^{\times}$, «~homothéties~».}

\page{60}
Notons que, si $Z_i$ \add{$\subset \mathrm{Gl}(V)$} est le groupe engendré
par $\tau_i$, on a $Z_i \cap k^{\times} = 1$ donc
$Z_i \hookrightarrow \mathrm{Aut}(C)$, donc $\nu_i$ prend une des valeurs
$2, 3, 4, 5$
\[
2 \leq \nu_0, \nu_1 \leq 5
\]
\struck{Il faut} Plus général, \struck{\ill{} \ill{}} \add{\ill{}} les
\uncertain{ps.\ réfl.}\ \add{strictes} de $G_0$, son \ill{}
\struck{\ill{} il \ill{} groupe engendré par} \add{\ill{}} ordre $\nu$ est
compris entre 2 et 5\ldots
\note{La lettre qu'il écrit pour le groupe engendré par $\tau_i$ a la forme
d'un 3 bouclé ; elle est rendue $Z_i$. Le reste de la page est blanc.}

\end{document}
